\documentclass[12pt]{article}

\usepackage{fontawesome}
\usepackage{hyperref}
\usepackage{xurl}
\usepackage[tagged, highstructure]{accessibility}
\usepackage{bookmark}
\usepackage{enumitem}

\hypersetup{
    colorlinks=false,
    pdfborder={0 0 0},
    pdftitle={Scripting with Variables},
    pdfauthor={Adrianna Holden-Gouveia},
    pdfsubject={Introduction to Linux - Lab Assignment},
    pdflang={en-US},
    pdfkeywords={lab assignment, Scripting with Variables}
}

% Configure PDF bookmarks for navigation
\bookmarksetup{
    numbered,
    open,
}

% Configure list spacing for better accessibility
\setlist{nosep}



\title{Scripting with Variables}
\author{
        Adrianna Holden-Gouveia \\
        Website: \href{https://aholdengouveia.name}{aholdengouveia.name}\\ 
        \date{\vspace{-5ex}}
        %Email: \href{mailto:admin@aholdengouveia.name}{admin@aholdengouveia.name} \\
%        \faLinkedin{: aholdengouveia} \\
        \faGithub {: aholdengouveia} \\
%        \faTwitter {: aholdengouveia} \\
        }

%S\date{\today}


\begin{document}    

\maketitle

\section*{Accessibility Notice}
This document is also available in HTML format at:

Link: \href{https://aholdengouveia.name/IntroLinux/labs/shellvar.html}{\textbf{https://aholdengouveia.name/IntroLinux/labs/shellvar.html}}

The HTML version provides enhanced accessibility features including keyboard navigation, screen reader support, responsive design, dark mode support, and high contrast options.


%\begin{abstract}
%This is a template for Linux Administration Lab work
%\end{abstract}
%\tableofcontents

\section*{Objectives:}
\begin{itemize}
    \item Learn how to do some more basic bash scripting. Become comfortable and competent with positional parameters, user input and backtics.
\end{itemize}
\section*{Complete the following problems}

References, a video, a PowerPoint and some notes are available at my website
Link: \href{https://www.aholdengouveia.name/IntroLinux/shellvar.html}{\textbf{https://www.aholdengouveia.name/IntroLinux/shellvar.html}}

\subsection*{How to turn in Scripts}
    \begin{enumerate}
        \item Testing must be submitted along with your code. You should submit a copy of the code and a screenshot or the results of your test so that I can see a successful run of your code. Every script should be tested.  A script that doesn't run is worth 0 points.
        \item Comments should be done on EVERY line of code, (this is very bad practice in the real world, never do this for a job) so that I can see that you understand the code that you are submitting. Every line should have an explanation of what that line does.  Comments are worth 5 points per script.
        \item For submission you need to submit a copy of your code (actual text NO screenshots of code allowed) and your successful test run of your code for every script.
\end{enumerate}

\subsection*{Scripts}
    \begin{enumerate}
        \item Write a script that will create a file with the date appended to the name. Date should come from the server not you.  The name of the file should be one of the towns in our state.
        \item Write a script that will ask the user for their name and then echo out a greeting with their name.  Include a comment with a hello and your name
        \item Write a script that will create a madlib about your favourite places. Link: \href{https://en.wikipedia.org/wiki/Mad_Libs}{\textbf{https://en.wikipedia.org/wiki/Mad_Libs}}  The user should be asked for the words that you want to insert, you must come up with the paragraph. The paragraph needs to be at least 3 lines.  The story should be about one of your favourite places to visit or somewhere in the world you'd want to live.  Make sure to include a code comment that includes one sentence on why you picked that location, be specific. Madlib shouldn't have more than 5 blanks to fill in
    \end{enumerate}



\section*{Deliverables}
One file that includes all the requested information with your name, the date, and the lab title, make sure to follow the requested pattern for scripts, an example is included with the first scripting lab.  You need 3 separate scripts for this lab
\end{document}